% ================================================================
% Capitolul 2 – Aspecte Teoretice
% Lucrare de licență: Kernel Trap – Sistem Inteligent de Prevenire
% a Intruziunilor cu Componente de Honeypot și Învățare Automată
% Autor: Banciu George-Horațiu
% Coordonator: Conf. Univ. Dr. Bufnea Darius Vasile
% ================================================================

\chapter{Aspecte teoretice}

Capitolul de față prezintă fundamentele conceptuale și tehnologice pe care se construiesc
sistemele moderne de detectare și prevenire a intruziunilor bazate pe monitorizarea la nivel
de kernel. Sunt tratate paradigmele IDS/IPS, mecanismele de interceptare a apelurilor de sistem
prin eBPF, instrumentele de colectare a evenimentelor, tehnicile de detecție a anomaliilor prin
învățare automată nesupervizată și principiile apărării active prin honeypot-uri.

% ----------------------------------------------------------------
\section{Sisteme de detectare și prevenire a intruziunilor}
% ----------------------------------------------------------------

Un \textit{sistem de detectare a intruziunilor} (\textit{Intrusion Detection System}, IDS) este
o componentă de securitate ce monitorizează în timp real activitatea unui sistem sau rețea în scopul
identificării comportamentelor suspecte sau a activităților neautorizate~\cite{liu2018hids}. Atunci
când un IDS este dotat și cu capacități de răspuns activ (blocarea conexiunilor, terminarea
proceselor sau redirecționarea traficului), acesta devine un \textit{sistem de prevenire a
intruziunilor} (\textit{Intrusion Prevention System}, IPS).

Din perspectiva punctului de colectare a datelor, sistemele IDS se clasifică în două mari categorii:

\begin{itemize}
  \item NIDS (\textit{Network-based IDS}): monitorizează traficul de rețea la nivel de
    pachete, identificând tipare de atac (scanări de porturi, flood-uri, exploit-uri cunoscute)
    fără acces la starea internă a gazdelor individuale.
  \item HIDS (\textit{Host-based IDS}): rulează direct pe gazda monitorizată și analizează
    evenimente produse la nivelul sistemului de operare (apeluri de sistem, jurnale de aplicație,
    accesuri la fișiere și modificări de configurație), oferind vizibilitate asupra comportamentului
    proceselor individuale~\cite{liu2018hids}.
\end{itemize}

Cele două abordări sunt complementare: NIDS detectează atacuri la nivel de transport, dar nu poate
inspecta comunicațiile criptate și nu are acces la comportamentul intern al proceselor; HIDS are
vizibilitate deplină asupra stării gazdei, dar nu observă activitatea laterală dintre alte mașini
din rețea.

Literatura de specialitate distinge două paradigme principale de detecție~\cite{liu2018hids,ring2021deep}:

\textbf{Detecția bazată pe semnături} (\textit{signature-based detection}) compară activitatea
observată cu o bază de cunoștințe de tipare de atac cunoscute. Rata de fals pozitive este scăzută
iar detecția este rapidă, dar metoda este prin definiție oarbă la amenințările noi (\textit{zero-day})
care nu figurează în baza de semnături. Soluții comerciale precum Snort sau Suricata se bazează
preponderent pe această abordare.

\textbf{Detecția bazată pe anomalii} (\textit{anomaly-based detection}) construiește un model al
comportamentului normal și semnalează orice deviere semnificativă de la acesta. Această paradigmă
poate detecta, în principiu, atacuri necunoscute anterior, deoarece nu depinde de semnăturile lor
specifice. Compromisul acceptat constă într-o rată mai ridicată de fals pozitive, generată de
variabilitatea naturală a comportamentului benign. Sistemele moderne combină adesea ambele abordări
pentru a valorifica avantajele fiecăreia~\cite{ring2021deep}.

% ----------------------------------------------------------------
\section{Apeluri de sistem și rolul lor în securitate}
% ----------------------------------------------------------------

Sistemul de operare Linux separă execuția în \textit{user space} (unde rulează aplicațiile
neprivilegiate) și \textit{kernel space}, unde se execută codul privilegiat al nucleului.
Această separare este impusă hardware prin inelele de protecție ale procesorului (Ring 0 pentru
kernel, Ring 3 pentru aplicații) și constituie fundația modelului de securitate al sistemului
de operare~\cite{bovet2005understanding}.

Comunicarea dintre cele două domenii se realizează exclusiv prin \textit{apeluri de sistem}
(\textit{syscalls}): puncte de intrare controlate prin care un program solicită nucleului operații
privilegiate. Nucleul Linux expune în prezent peste 300 de apeluri de sistem, grupate funcțional
în categorii: gestionarea proceselor (\texttt{fork}, \texttt{execve}, \texttt{exit}), accesul la
fișiere (\texttt{open}, \texttt{read}, \texttt{write}, \texttt{unlink}), comunicația prin rețea
(\texttt{socket}, \texttt{connect}, \texttt{bind}), gestionarea utilizatorilor și privilegiilor
(\texttt{setuid}, \texttt{setgid}, \texttt{capset}) și controlul proceselor
(\texttt{kill}, \texttt{ptrace})~\cite{kerrisk2010linux}.

Orice acțiune relevantă din perspectiva securității trece obligatoriu prin interfața syscall.
Un atacator care rulează un exploit, escaladează privilegii sau exfiltrează date nu poate ocoli
această interfață, indiferent de tehnica folosită. Această proprietate face din secvențele de
syscall-uri o sursă ideală de telemetrie pentru sistemele HIDS.

Spre deosebire de logurile aplicative (care pot fi falsificate sau suprimate de un atacator cu
privilegii ridicate), evenimentele de la nivelul syscall-urilor sunt inspecționate de nucleu
înainte ca aplicația să aibă posibilitatea de a interveni~\cite{tracee2024}. Astfel, chiar dacă
un atacator compromite complet o aplicație, comportamentul acesteia rămâne vizibil prin
instrumentarea nucleului.

Din perspectiva unui HIDS, secvențele de syscall-uri permit identificarea unor tipare
comportamentale caracteristice atacurilor: un \texttt{execve} urmat imediat de \texttt{setuid}
și \texttt{capset} sugerează o tentativă de escaladare de privilegii; o explozie de apeluri
\texttt{connect} către adrese IP externe indică o potențială activitate de tip Command and Control;
o secvență \texttt{open}/\texttt{read} pe fișierele \texttt{/etc/shadow} sau \texttt{/etc/passwd}
semnalează recunoaștere locală a sistemului~\cite{liu2018hids}.

% ----------------------------------------------------------------
\section{Extended Berkeley Packet Filter (eBPF)}
% ----------------------------------------------------------------

Berkeley Packet Filter (BPF), introdus în 1992 de McCanne și Jacobson ca mecanism de filtrare
a pachetelor de rețea la nivel de kernel~\cite{mccanne1993bsd}, a fost extins radical începând
cu nucleul Linux 3.18 (2014) sub forma \textit{extended BPF} (eBPF)~\cite{gregg2019bpf}. Dacă
BPF-ul clasic opera exclusiv pe pachete de rețea, eBPF permite atașarea de programe la orice
punct de instrumentare din nucleu: syscall-uri, kprobes, uprobes, tracepoints și evenimente LSM
(\textit{Linux Security Modules}).

Arhitectural, eBPF funcționează ca o mașină virtuală ușoară integrată în nucleu, cu un set de
instrucțiuni propriu (64 de biți, 11 regiștri generali) și un model de execuție bazat pe evenimente.
Programele sunt scrise în C restricționat, compilate în bytecode și supuse unui \textit{verificator
static} care garantează absența buclelor infinite, a acceselor invalide la memorie și a scurgerilor
de resurse~\cite{gregg2019bpf}. Verificatorul analizează toate căile posibile de execuție
ale programului înainte de încărcare, respingând programele care nu pot fi demonstrate sigure.
Bytecode-ul validat este ulterior compilat JIT (\textit{Just-in-Time}) în instrucțiuni native ale
procesorului.

Starea este partajată între programe și user space prin \textit{eBPF Maps}, structuri cheie-valoare
rezidente în kernel space, accesibile din ambele domenii. Apelurile din programele eBPF către
funcțiile nucleului se fac exclusiv prin \textit{helper functions} cu interfață stabilă, garantând
compatibilitatea între versiuni diferite de kernel. Figura~\ref{fig:cap2-ebpf} sintetizează acest
flux, de la apelul de sistem produs în spațiul utilizator până la citirea evenimentului filtrat,
într-un map, de către un agent din spațiul utilizator precum Tracee.

\begin{figure}[h]
  \centering
  \begin{tikzpicture}[
      font=\small,
      box/.style={rectangle, rounded corners=3pt, draw, line width=0.4pt,
        minimum width=3cm, minimum height=1cm, align=center, inner sep=3pt},
      flow/.style={-{Stealth[length=2.2mm]}, line width=0.5pt, gray!70!black},
      lbl/.style={font=\scriptsize\itshape, text=black!65},
    ]
    \definecolor{cap2ssh}{RGB}{230,241,251}
    \definecolor{cap2tracee}{RGB}{225,245,238}
    \definecolor{cap2gray}{RGB}{241,239,232}
    \definecolor{cap2prog}{RGB}{238,237,254}

    \draw[dashed, gray!60, line width=0.4pt] (-2,-0.75) rectangle (11,1.05);
    \node[anchor=north west, font=\scriptsize, text=gray!70!black] at (-1.95,1.02) {spațiul utilizator};
    \draw[dashed, gray!60, line width=0.4pt] (-2,-3.85) rectangle (11,-2.2);
    \node[anchor=north west, font=\scriptsize, text=gray!70!black] at (-1.95,-2.25) {spațiul nucleu};
    \draw[dashed, gray!55, line width=0.4pt] (-2,-1.5) -- (11,-1.5);
    \node[font=\scriptsize\itshape, text=black!60] at (4.5,-1.3) {interfața de apeluri de sistem};

    \node[box, fill=cap2ssh]    (ssh)    at (0,0)    {\textbf{Sesiune SSH}\\[1pt]\scriptsize proces utilizator};
    \node[box, fill=cap2tracee] (tracee) at (9,0)    {\textbf{Tracee}\\[1pt]\scriptsize agent user-space};
    \node[box, fill=cap2gray]   (tp)     at (0,-3)   {\textbf{Tracepoint}\\[1pt]\scriptsize hook syscall};
    \node[box, fill=cap2prog]   (prog)   at (4.5,-3) {\textbf{Program eBPF}\\[1pt]\scriptsize filtrare în kernel};
    \node[box, fill=cap2gray]   (map)    at (9,-3)   {\textbf{Map eBPF}\\[1pt]\scriptsize tampon partajat};

    \draw[flow] (ssh.south) -- node[lbl, left] {apel de sistem} (tp.north);
    \draw[flow] (tp.east) -- (prog.west);
    \draw[flow] (prog.east) -- (map.west);
    \draw[flow] (map.north) -- node[lbl, right] {citire map} (tracee.south);
  \end{tikzpicture}
  \caption[Fluxul de colectare prin eBPF]{Fluxul de colectare prin eBPF. Un apel de sistem produs de o sesiune SSH traversează
  interfața de apeluri de sistem; în nucleu, un program eBPF atașat la un tracepoint filtrează
  evenimentul și îl scrie într-un map, de unde agentul Tracee îl citește în spațiul utilizator.}
  \label{fig:cap2-ebpf}
\end{figure}

Înainte de eBPF, instrumentarea nucleului la runtime necesita una din două abordări, ambele
problematice în producție: modificarea codului sursă al nucleului (imposibilă fără acces la sursă
și fără repornire) sau încărcarea de module de nucleu (\textit{Loadable Kernel Modules}, LKM).
Modulele de nucleu rulează în Ring 0, fără niciun mecanism de verificare prealabilă, iar un bug
poate produce panică de kernel sau compromiterea întregului sistem.

eBPF elimină aceste limitări: programele se încarcă la runtime, fără repornire, și sunt garantate
sigure de verificator înainte de execuție. Față de mecanisme alternative precum \texttt{auditd}
sau \texttt{ptrace}, eBPF prezintă cost suplimentar semnificativ mai redus, deoarece filtrarea și
agregarea evenimentelor se realizează direct în kernel space, reducând la minimum cantitatea de
date transferate către user space~\cite{gregg2019bpf}.

% ----------------------------------------------------------------
\section{Tracee: instrumentul de colectare a evenimentelor}
% ----------------------------------------------------------------

Tracee este un instrument open-source de securitate runtime pentru Linux, dezvoltat și menținut
de Aqua Security~\cite{aqua2024tracee}. Folosește eBPF pentru a intercepta evenimentele de sistem
la runtime și a le expune ca fluxuri de date structurate în format JSON, acoperind un spectru larg
de activitate: crearea și terminarea de procese, accesul la fișiere, conexiunile de rețea,
modificările de permisiuni și evenimentele LSM.

Tracee este compus din două componente principale: \textbf{Tracee-eBPF}, care compilează și
încarcă programele eBPF în nucleu și expune fluxul brut de evenimente, și \textbf{Tracee-Rules},
un motor de detectare bazat pe reguli scrise în Go sau Rego (limbajul de politici al Open Policy
Agent). O caracteristică notabilă este preferința pentru evenimentele
LSM față de syscall-urile brute (de exemplu, \texttt{security\_file\_open} în locul \texttt{openat}),
care oferă informații mai complete despre contextul de securitate și sunt mai puțin susceptibile
la atacuri de tip TOCTOU (\textit{Time-of-Check to Time-of-Use})~\cite{aqua2024tracee}.

Portabilitatea pe kernel-uri diferite este asigurată prin suportul \textit{Compile Once – Run
Everywhere} (CO:RE), bazat pe \textit{BPF Type Format} (BTF): pe nucleele cu BTF activat, nu
sunt necesare nici recompilarea probe-urilor, nici prezența header-elor de nucleu pe sistemul
țintă. Aceasta face din Tracee o soluție practică pentru deployment pe
infrastructuri eterogene.

% ----------------------------------------------------------------
\section{Honeypot-uri și apărarea activă}
% ----------------------------------------------------------------

Un \textit{honeypot} este un sistem de securitate proiectat deliberat să fie compromis, cu scopul
de a atrage atacatorii, de a le studia tehnicile și de a le deturnă resursele și atenția de la
sistemele reale~\cite{spitzner2002honeypots}. Spre deosebire de sistemele de producție, un honeypot
nu conține date sau servicii reale, iar orice interacțiune cu el este, prin definiție, suspectă.

Taxonomia standard clasifică honeypot-urile după două criterii principale:

\textbf{Nivelul de interacțiune} distinge între:
\begin{itemize}
  \item \textit{Honeypot-uri cu interacțiune redusă} (\textit{low-interaction}): emulează servicii
    și sisteme de operare fără a rula un sistem real. Sunt ușor de implementat și cu risc minim
    de compromitere, dar oferă informații limitate despre tehnicile atacatorului. Exemple tipice:
    Honeyd, OpenCanary.
  \item \textit{Honeypot-uri cu interacțiune ridicată} (\textit{high-interaction}): expun sisteme
    reale (sau mașini virtuale complete) cu servicii funcționale. Colectează informații detaliate
    despre comportamentul atacatorului, inclusiv tool-uri folosite, comenzi executate și ținte
    vizate ulterior, dar prezintă risc mai ridicat dacă atacatorul reușește să evadeze din
    mediul controlat~\cite{spitzner2002honeypots}.
\end{itemize}

\textbf{Scopul de utilizare} distinge între:
\begin{itemize}
  \item \textit{Honeypot-uri de cercetare}: implementate de organizații de securitate pentru
    studierea tehnicilor de atac emergente, colectarea de mostre de malware și înțelegerea
    comportamentului atacatorilor.
  \item \textit{Honeypot-uri de producție}: integrate în rețelele organizațiilor pentru a detecta
    activitate neautorizată internă și a alerta echipele de securitate.
\end{itemize}

Honeypot-urile tradiționale funcționează pasiv: atacatorul trebuie să ajungă la capcană din
proprie inițiativă, alegând să interacționeze cu un serviciu-momeală. Această abordare prezintă
o limitare fundamentală în scenariile în care atacatorul a compromis deja un sistem legitim,
deoarece el nu are niciun motiv să migreze voluntar către un honeypot.

O abordare mai avansată, denumită \textit{apărare activă} (\textit{active deception}), inversează
această logică prin redirecționarea transparentă a unui atacator detectat de pe sistemul compromis
către un honeypot, fără ca acesta să conștientizeze tranziția. Mecanismul se bazează tipic pe
reguli NAT (\texttt{iptables}/\texttt{nftables}), care interceptează traficul la nivel de rețea
și îl deviază înainte ca pachetele să ajungă la destinația originală~\cite{spitzner2002honeypots}.
Avantajul față de abordarea pasivă constă în faptul că atacatorul nu trebuie să aleagă
honeypot-ul din proprie inițiativă, ci este direcționat acolo în momentul în care comportamentul
său declanșează un semnal de alertă, indiferent că a compromis deja un serviciu real.

Un concept complementar honeypot-urilor îl reprezintă \textit{honeytoken-ul}: un artefact digital
fals (fișier, cheie de acces, înregistrare în bază de date) plasat intenționat într-un sistem
real sau într-un honeypot, proiectat să nu fie niciodată accesat de utilizatori legitimi. Orice
interacțiune cu un honeytoken constituie, prin definiție, un indicator incontestabil de
compromitere~\cite{spitzner2002honeypots}. Spre deosebire de detecția bazată pe anomalii
statistice, care produce un scor continuu cu risc de fals pozitive, accesul la un honeytoken
este un semnal binar și precis: niciun utilizator legitim nu are motiv să deschidă un fișier de
credențiale pe care el însuși nu l-a creat. Honeytokens-urile pot lua diverse forme: fișiere cu
credențiale fictive, chei SSH false, token-uri de API nevalide sau înregistrări-capcană în baze
de date. Sunt utilizate atât în medii de producție pentru detectarea accesului neautorizat,
cât și în honeypot-uri pentru a monitoriza comportamentul post-compromitere al atacatorilor.

% ----------------------------------------------------------------
\section{Algoritmul Isolation Forest}
% ----------------------------------------------------------------

Isolation Forest (\textit{iForest}), propus de Liu, Ting și Zhou în 2008~\cite{liu2008iforest},
adoptă o paradigmă diferită față de metodele clasice de detecție a anomaliilor (LOF, k-NN,
One-Class SVM): în loc să modeleze distribuția normală și să identifice instanțele care se
abat de la ea, izolează explicit anomaliile prin partiționări recursive aleatorii.

Intuiția algoritmului se bazează pe două proprietăți ale anomaliilor: acestea sunt \textit{rare}
(puține la număr față de instanțele normale) și \textit{diferite} (ocupă regiuni izolate ale
spațiului de caracteristici). Ca urmare, ele pot fi separate de restul datelor printr-un număr
semnificativ mai mic de partiționări aleatorii decât instanțele normale, care sunt dense și
necesită multe partiționări pentru a fi izolate~\cite{liu2012isolation}
(Figura~\ref{fig:cap2-iforest}).

\begin{figure}[h]
  \centering
  \begin{tikzpicture}[
      font=\small,
      boxc/.style={rectangle, rounded corners=3pt, draw, line width=0.4pt,
        minimum width=4.4cm, minimum height=1.1cm, align=center, inner sep=3pt},
      flow/.style={-{Stealth[length=2.2mm]}, line width=0.5pt, gray!70!black},
    ]
    \definecolor{cap2anom}{RGB}{226,75,74}
    \definecolor{cap2norm}{RGB}{55,138,221}
    \definecolor{cap2redbox}{RGB}{252,235,235}
    \definecolor{cap2bluebox}{RGB}{230,241,251}

    \node[anchor=south west, font=\scriptsize, text=black!65] at (0,4.55) {spațiul caracteristicilor};
    \draw[line width=0.4pt, black!55] (0,0) rectangle (5,4.5);

    \draw[gray!55, line width=0.3pt] (2.2,0) -- (2.2,4.5);
    \draw[gray!55, line width=0.3pt] (0,3.0) -- (2.2,3.0);
    \draw[gray!55, line width=0.3pt] (3.5,0) -- (3.5,4.5);
    \draw[gray!55, line width=0.3pt] (2.2,2.0) -- (5,2.0);
    \draw[gray!55, line width=0.3pt] (4.2,2.0) -- (4.2,4.5);
    \draw[gray!55, line width=0.3pt] (2.85,0) -- (2.85,2.0);

    \draw[cap2anom, dashed, line width=0.5pt] (0,3.0) rectangle (2.2,4.5);
    \fill[cap2anom] (1.0,3.6) circle (0.09);
    \node[font=\scriptsize, text=cap2anom] at (1.0,4.15) {anomalie};

    \foreach \p in {(3.0,1.5),(3.4,1.3),(3.7,1.6),(3.2,1.0),(3.6,0.9),(3.9,1.3),(3.3,1.8),(3.0,1.0),(4.0,1.7),(4.2,1.0),(2.9,1.7),(4.1,0.8)}
      \fill[cap2norm] \p circle (0.08);
    \node[font=\scriptsize, text=cap2norm!70!black] at (3.5,0.35) {puncte normale};

    \node[boxc, fill=cap2redbox,  draw=cap2anom]      (short) at (8.7,3.4) {\textbf{Cale scurtă de izolare}\\[1pt]\scriptsize puține tăieturi = anomalie};
    \node[boxc, fill=cap2bluebox, draw=blue!45!black] (long)  at (8.7,1.2) {\textbf{Cale lungă de izolare}\\[1pt]\scriptsize multe tăieturi = normal};

    \draw[flow] (5,3.6) -- (short.west);
    \draw[flow] (5,1.4) -- (long.west);
  \end{tikzpicture}
  \caption[Intuiția algoritmului Isolation Forest]{Intuiția algoritmului Isolation Forest. O anomalie, situată într-o regiune rară a
  spațiului de caracteristici, este izolată din puține tăieturi aleatorii (cale scurtă în arbore),
  în timp ce un punct dintr-un cluster dens necesită multe tăieturi (cale lungă) pentru a fi
  separat.}
  \label{fig:cap2-iforest}
\end{figure}

Algoritmul construiește o pădure de \textit{isolation trees} (iTrees). Fiecare iTree este un
arbore binar construit prin selecția aleatorie repetată a unui atribut $q$ și a unui prag de
partiționare $p$ ales uniform în intervalul de valori al atributului, până când fiecare instanță
este izolată într-un nod frunză. O pădure constă din $t$ astfel de arbori, fiecare antrenat
pe un subsample de dimensiune $\psi$; valorile implicite sunt $t = 100$, $\psi = 256$~\cite{liu2008iforest}.

Lungimea căii $h(x)$ printr-un iTree reprezintă numărul de muchii parcurse de la rădăcină până
la nodul de izolare al instanței $x$. Anomaliile, mai ușor de izolat, prezintă căi mai scurte;
instanțele normale, mai greu de separat, prezintă căi mai lungi. Scorul de anomalie normalizat
este definit ca:
\[
    s(x, \psi) = 2^{-\dfrac{E[h(x)]}{c(\psi)}}
\]
unde $E[h(x)]$ este media lungimilor de cale pe toți $t$ arborii, iar $c(\psi)$ este lungimea
medie de cale așteptată pentru un arbore binar de căutare construit pe $\psi$
instanțe~\cite{liu2012isolation}. Un scor $s \approx 1$ indică anomalie certă, $s \approx 0.5$
instanță nedistinctivă, iar $s \approx 0$ comportament cu certitudine normal.

Față de metodele bazate pe distanță sau densitate, iForest are complexitate liniară $O(n)$,
memorie redusă prin subsampling fix și nu necesită date etichetate cu atacuri, proprietăți
esențiale într-un context de securitate unde datele de atac sunt rare sau inexistente la
momentul antrenării~\cite{liu2008iforest,liu2012isolation}.

% ----------------------------------------------------------------
\section{Setul de date BETH}
% ----------------------------------------------------------------

Setul de date BETH (\textit{BPF-Extended Tracking Honeypot}), prezentat de Highnam et al. la
ICML Workshop on Uncertainty and Robustness in Deep Learning (2021)~\cite{highnam2021beth}, este
unul dintre puținele seturi de date publice de evenimente de syscall colectate dintr-un scenariu
real de honeypot în mediu cloud. Conține peste 8 milioane de evenimente capturate de pe 23 de
gazde honeypot expuse pe internet timp de mai multe zile.

Fiecare eveniment este reprezentat prin 14 caracteristici brute: \texttt{processId},
\texttt{parentProcessId}, \texttt{userId}, \texttt{mountNamespace}, \texttt{eventId},
\texttt{argsNum}, \texttt{returnValue} și altele, și este etichetat binar: benign (\texttt{evil=0})
sau malițios (\texttt{evil=1}). Setul este partiționat 60/20/20 (antrenament/validare/test);
important, datele de atac apar \textit{exclusiv} în partiția de test, reflectând paradigma
realistă în care un model este antrenat pe comportament benign și evaluat ulterior pe atacuri.
Atacul reprezentat constă în instanțierea unui nod de botnet: exploatare inițială, escaladare
de privilegii și activitate de Command and Control (C\&C).

Față de alte seturi de date sintetice sau de laborator, BETH prezintă avantajul că reflectă
comportamentul real al unui atacator pe servere Linux expuse la internet, colectat prin aceeași
tehnologie eBPF utilizată de instrumentele moderne de monitorizare. Isolation Forest atinge
AUROC de 0.850 pe acest set, cel mai bun rezultat dintre modelele clasice nesupervizate testate
în literatura de specialitate~\cite{highnam2021beth}.

% ----------------------------------------------------------------
\section{Redis Streams}
% ----------------------------------------------------------------

Redis Streams, introdusă în Redis 5.0 (2018), este o structură de date de tip jurnal
\textit{append-only} cu acces aleator în timp $O(1)$, destinată scenariilor de mesagerie și
procesare de evenimente în timp real~\cite{redis2024streams}. Fiecare intrare
(\textit{entry}) dintr-un stream este identificată printr-un ID monoton crescător (compus din
timestamp în milisecunde și un număr de secvență) și conține un set arbitrar de perechi
câmp-valoare.

Spre deosebire de Redis Pub/Sub,care pierde mesajele dacă niciun subscriber nu este activ
în momentul publicării, Streams persistă intrările pe disc și le face disponibile pentru citire ulterioară. Mecanismul de \textit{consumer groups}
permite mai multor consumatori să citească din același stream în paralel, fiecare primind un
subset disjunct de mesaje; un mesaj rămâne în starea \textit{pending} până la confirmarea
explicită prin comanda \texttt{XACK}, garantând că niciun eveniment nu se pierde în caz de
cădere a consumatorului.

Aceste caracteristici fac din Redis Streams o soluție potrivită pentru arhitecturi de tip
pipeline în care un producător de evenimente (agentul de monitorizare) trebuie decuplat de
consumatorul de analiză, cu garanții de livrare și fără introducerea complexității operaționale
a unui broker distribuit dedicat.

% ----------------------------------------------------------------
\section{Containerizarea Docker și rolul în mecanismul honeypot}
% ----------------------------------------------------------------

Docker permite ambalarea unei aplicații împreună cu toate dependențele sale într-o unitate
portabilă și reproductibilă numită \textit{container}~\cite{merkel2014docker}. Spre deosebire
de mașinile virtuale, care virtualizează hardware-ul complet și rulează un nucleu propriu,
containerele partajează nucleul sistemului gazdă și se izolează exclusiv prin mecanisme
software ale nucleului Linux:

\begin{itemize}
  \item \textit{Namespaces}: creează viziuni izolate ale resurselor sistemului pentru fiecare
    container: spații separate de PID (procesele unui container nu pot vedea procesele altui
    container), rețea (interfețe și tabele de rutare proprii), sistem de fișiere
    (\textit{mount namespace}), utilizatori (\textit{user namespace}) și hostname.
  \item \textit{Control groups} (\textit{cgroups}): limitează cantitatea de resurse (CPU,
    memorie, I/O pe disc, lățime de bandă) pe care un container o poate consuma, prevenind
    consumul abuziv al resurselor gazdei.
\end{itemize}

Izolarea prin containere se dovedește utilă și în contextul apărării active: un container
honeypot cu interacțiune ridicată poate expune un shell real și unelte uzuale, simulând un
sistem Linux funcțional, fără ca activitatea atacatorului să afecteze infrastructura de
producție. Tranzițional, regulile NAT (\texttt{iptables}, tabela \texttt{nat}, lanțul
\texttt{PREROUTING}) pot intercepta traficul destinat unui serviciu legitim și îl pot devia
transparent către portul corespunzător al containerului honeypot, înainte ca pachetele să
ajungă la destinația originală. Din perspectiva atacatorului, sesiunea continuă fără întrerupere
vizibilă, dar el interacționează cu un mediu controlat și monitorizat, izolat complet de
sistemele de producție~\cite{spitzner2002honeypots,kerrisk2010linux}.
